% =====================================================================
\section{主定理与其地位}\label{sec:dy-theorem}
% =====================================================================

现在可以把前几节的结果装配成单举 DY 的因子化定理。装配本身是一个演绎：把模板定理\cref{thm:method-template}的假设 (H1)--(H5)（该定理本身是全阶论证，\tagA）逐条对应到本章的陈述。这里的难点不在演绎，而在于如实交代每一条假设的地位。我们按四个层次给出这一交代：文献中的论证；把它写成的精确条件定理链；以标签列出的、尚未证明的输入；以及反例的情况。

\begin{theorem}[单举 DY 的共线因子化 \tagA]\label{thm:dy-factorization}
对单举过程 $A+B\to\gamma^*+X$（$Q\gg\Lqcd$，$\tau$ 与 $y$ 固定），
\begin{equation}\label{eq:dy-factorization}
\frac{\dd\sigma}{\dd Q^2\,\dd y}=\sum_{i,j}\int_{x_A}^1\frac{\dd\xi_A}{\xi_A}\int_{x_B}^1\frac{\dd\xi_B}{\xi_B}\,f_{i/A}(\xi_A;\mu)\,f_{j/B}(\xi_B;\mu)\,\mathcal H_{ij}\Big(\frac{x_A}{\xi_A},\frac{x_B}{\xi_B},\frac Q\mu,\as(\mu)\Big)+\order{\Big(\frac{\Lambda}{Q}\Big)^p},
\end{equation}
其中 $f_{i/A}$、$f_{j/B}$ 是\cref{def:pdf-quark,def:pdf-gluon}中的 $\MSbar$ 部分子分布，与 DIS 中的是同一组函数；$\mathcal H_{ij}$ 与外态无关，可用部分子态逐阶计算；余项的幂次 $p>0$ 不由证明确定（\tagC，预期 $p=2$）。\\
非微扰地位：\cref{asm:intro-bridge}（MA1、MA2、MA4）。
% errtype: E-S
\uses{thm:method-template,prop:dy-regions,lem:dy-fsi,prop:dy-deformation,lem:dy-soft,prop:dy-universality,thm:appO-main,cor:appO-interface,prop:appN-dy,def:pdf-quark,def:pdf-gluon,asm:intro-bridge,asm:appO-lcopt,asm:appO-envelope,asm:appO-mixing,asm:appO-deformation,asm:appO-rewritten-cs,asm:appN-collins,asm:appN-cutsum}
\end{theorem}

\begin{proofsketch}
由\cref{thm:method-template}（全阶论证，\tagA），只需对 DY 验证其假设 (H1)--(H5)，并把各步的精确形式接到附录的条件定理上。
(H1) 领头区域由\cref{prop:dy-regions}给出，其中包含旁观者之间的 Glauber 子区域；积分范数下的幂次计数是模板定理 (H1) 的前提。
(H2) $H$ 与 $J$ 的接口处的共线近似与纵向胶子的退耦同 DIS：$A$ 共线的纵向胶子附着在 $H$ 或 $J_B$ 上，在 $A$ 的共线因子中给出指向过去的 $W[\cdot,-\infty n_B]$；由\cref{prop:dy-universality}，它与 DIS 中的 $W[\cdot,\infty n_B]$ 给出同一个算符。
(H3) Glauber 区：割图求和之后把软动量沿 $e_G$ 形变出 Glauber 区，软胶子退耦为指向过去的 Wilson 线（\cref{lem:dy-fsi,prop:dy-deformation}）。
(H4) 不受测量约束的末态相互作用由幺正性相消，对 $q_T$ 积分的截面软因子为 $1$（\cref{lem:dy-fsi,lem:dy-soft}）。
(H5) 减除项与部分子分布算符的重整化矩阵元一致，同 DIS；DY 所需的积分形式是\cref{prop:appN-dy}：在\cref{asm:appN-collins}、\cref{asm:appN-cutsum}、\cref{asm:appO-rewritten-cs}、\cref{asm:appO-deformation}(b)(c) 之下，$\big|\int F_G-\sum_R\int C_RF_G\big|\le C\lambda^aL^K$，并且在每个由 $R_0$ 主导的块上，减除项可以用单侧改写的形式计算（\cref{cor:appO-interface}(e)），其中 Wilson 线指向过去。
把 (H1)--(H5) 代入模板定理，软因子为 $1$，得 $\sigma=\sum_{ij}\mathcal H_{ij}\otimes f_{i/A}\otimes f_{j/B}$ 加上被幂次压低的余项；$\mathcal H_{ij}$ 是部分子截面减去部分子分布的卷积（质量因子化），因而与外态无关。各步的精确条件形式依次是\cref{thm:appO-main}、\cref{cor:appO-interface}、\cref{prop:appN-dy}（均为 \tagR，$R_{\mathrm{cond}}$），其前提为\cref{asm:appO-lcopt}、\cref{asm:appO-envelope}、\cref{asm:appO-mixing}、\cref{asm:appO-deformation}、\cref{asm:appO-rewritten-cs}以及上面的附录~N 输入。完整的论证见\cite[第~14.4~节]{Collins2011}；精确形式见\cref{sec:appO-main}。
\end{proofsketch}

\paragraph{地位与缺口}\tagA。我们分四层说明。

\emph{(1) 文献中的论证。}CSS（1985、1988）与 Bodwin（1985）给出了全阶论证的关键步骤\cite{CollinsSoperSterman1985b,CollinsSoperSterman1988,Bodwin1985}：Glauber 区的围道变形、末态相互作用的幺正性相消、软因子为 $1$；两圈上旁观者的非因子化项之和为零\cite{Lindsay:1982pu,Lindsay:1982gd}，程函“介子”模型中的两圈相消见\cite{Brandt:1988xt}（均据其摘要）。Collins 的教科书说明，其第 14.4 节的处理遵循 CSS 1988，并作了改进、改正了其中的错误（他没有具体列出这些错误）\cite[第~14.4~节，第~553~页]{Collins2011}。他在第 14.4.8 节“Issues remaining”中列出：(i) 证明依赖微扰论的动量分类，非微扰的“分级快度”Glauber 交换（多周边、类 Regge）没有被覆盖；(ii) Feynman 规范下领头区域不完备，额外的不连通硬散射须对图求和后才相消\cite{Labastida:1984gy}；(iii) “The current versions of genuine proofs are evidently formidable”（现有的真正证明显然是令人生畏的）\cite[第~14.4.8~节]{Collins2011}。Rothstein 与 Stewart 的 Glauber 算符框架复现了旁观者相位的相消，但他们自己说明这一相消“taken alone it is far from a proof of factorization, even in the abelian case”（单独来看远不是因子化的证明，即使在阿贝尔情形）\cite[第~11.1~节]{RothsteinStewart2016}；据我们截至 2026 年 10 月在 INSPIRE 上的文献检索，他们宣布的完整 SCET 证明尚未发表。

\begin{table}[!tbp]
\centering
\caption{DY 证明的逐步状态（附录~O 的步骤 (O-1)--(O-8)，据表~\ref{tab:appO-checklist}）}
\label{tab:dy-status}
\footnotesize
\setlength{\tabcolsep}{3pt}
\renewcommand{\arraystretch}{1.18}
\begin{tabularx}{\textwidth}{@{}>{\raggedright\arraybackslash}p{1.95cm}>{\raggedright\arraybackslash}X>{\raggedright\arraybackslash}p{2.1cm}>{\raggedright\arraybackslash}p{1.25cm}>{\raggedright\arraybackslash}p{2.75cm}>{\raggedright\arraybackslash}p{2.0cm}@{}}
\toprule
步骤 & 陈述（标签） & 地位 & 误差类型 & 所用假设 & 文献\\
\midrule
(O-1) 区域、定义域与窗口 & \cref{prop:dy-regions}；\cref{def:appO-leading-data}、\cref{def:appO-windows}、\cref{lem:appO-glauber-pc} & 定义；幂次计数为标度陈述 & E-S；E-P & \cref{asm:appN-envelope}、\cref{asm:appN-separation} & \cite{CollinsSoperSterman1985b}；\cite[14.3~节]{Collins2011}\\
(O-2) 光锥序中的领头幂次形式 & \cref{lem:appO-lcopt}、\cref{lem:appO-hard-locality}、\cref{lem:dy-poles-orientation} & 在假设下完整 & 精确 & \cref{asm:appO-lcopt} & \cite{Kogut:1969xa,Chang:1973qi,Ligterink:1994tm,AybatSterman2009}\\
(O-3) 割图求和与单侧改写 & \cref{lem:dy-fsi}（\cref{lem:appO-unitarity}）、\cref{lem:appO-equal-time}、\cref{lem:appO-joint-localization}、\cref{prop:appO-ss-cancel}、\cref{prop:appO-one-sided} & 在假设下完整 & 精确；E-F；E-W & \cref{asm:appO-lcopt}、\cref{asm:appO-envelope}、\cref{asm:appO-mixing} & \cite[14.4.1~节]{Collins2011}；\cite[4~节]{Diehl:2015bca}\\
(O-4) 减除项与改写相容 & \cref{prop:appO-one-sided}（链项） & 在假设下完整 & E-W & \cref{asm:appN-separation}(ii)、\cref{asm:appO-envelope} & ---\\
(O-5) 形变目标 & \cref{prop:dy-deformation}；\cref{lem:appO-no-crossing}、\cref{prop:appO-target} & 对 $\mathcal G_0$ 与分级构型完整，一般情形为假设 & 精确（几何） & \cref{asm:appO-deformation}(a)--(d) & \cite{CollinsSoperSterman1988,AybatSterman2009}\\
(O-6) 程函近似与指向过去的 Wilson 线 & \cref{prop:appO-eikonal}、\cref{prop:appO-past-wilson} & 程函界完整；Wilson 线为 \tagA & E-P & \cref{asm:appO-rewritten-cs}、\cref{asm:appN-ward-product} & \cite{GrammerYennie1973}；\cite[14.4~节]{Collins2011}\\
(O-7) 软因子为 1；共线因子为 DIS 的 PDF & \cref{lem:dy-soft}（\cref{lem:method-soft-unitarity}）、\cref{prop:dy-universality} & 完整（$R_{\mathrm{def}}$） & 精确 & --- & \cite[14.4.6~节]{Collins2011}\\
(O-8) 与附录~N 的接口 & \cref{thm:appO-main}、\cref{cor:appO-interface}、\cref{prop:appN-dy} & 在假设下完整（$R_{\mathrm{cond}}$） & E-W$\to$E-N & 以上全部，及\cref{asm:appN-collins}、\cref{asm:appN-cutsum} & ---\\
\bottomrule
\end{tabularx}
\end{table}

\emph{(2) 精确形式。}\cref{app:appO}把 (H3)(H4) 写成条件定理\cref{thm:appO-main}（\tagR，$R_{\mathrm{cond}}$），经接口\cref{cor:appO-interface}接到\cref{app:appN}的\cref{prop:appN-dy}（\tagR，$R_{\mathrm{cond}}$），由此得到因子化形式。因此，DY 的共线因子化等于“割图求和后的 Glauber 相消”（\cref{app:appO}）加上减除方法中组合与幂次计数的部分，而后者与 $e^+e^-$、DIS 中的完全相同。这条链的内容等级是 \tagA。

\emph{(3) 以标签列出的开放输入。}\cref{asm:appO-lcopt}（光锥序微扰论的适用性）；\cref{asm:appO-envelope}（纤维积分误差的可积包络）；\cref{asm:appO-mixing}(c)（重合连接的完整图集相消；对程函线之间的双胶子交换在 $\order{g^4}$ 成立，见\cref{prop:dy-eikonal-twogluon}；对复合旁观者的未割线，在全局两时模型中有模型证据，见\cref{prop:appO-twoexch-spectator}）；\cref{asm:appO-deformation}(b)(c)(d)（\emph{即使对阿贝尔图也需要}）；\cref{asm:appO-rewritten-cs}；附录~N 的一般输入\cref{asm:appN-collins}、\cref{asm:appN-cutsum}，以及经\cref{thm:method-template}进入的\cref{asm:appN-envelope}、\cref{asm:appN-separation}、\cref{asm:appN-ward-product}、\cref{asm:appN-values}、\cref{asm:appN-uv}、\cref{asm:appN-multren}；非微扰的桥接\cref{asm:intro-bridge}。此外，该论证中被明确假定的一步（\cref{asm:appN-brst-nonlocal}，\tagC，多定域光线算符的 BRST 分类）是\cref{thm:method-template}中点名的具名缺口，在此重申；它不是本定理的前提。除约定\cref{asm:appN-values}之外，全部技术性输入都是 \tagA，所以本定理是 \tagA。各步的误差类型见\cref{tab:dy-status}：标度陈述（E-S）、逐点界（E-P）、纤维积分界（E-F）与弱等价（E-W）。

\emph{(4) 反例。}没有已知的领头幂次反例；文献中“严格窗口”的非幺正残差是截断的假象（\cref{prop:dy-eikonal-twogluon}）。混合光锥序与“钳夹变为初、末态形变方向的不匹配”的看法应归功于 Aybat 与 Sterman\cite{AybatSterman2009}（这是一个论证纲要，不是证明）；Diehl 等人给出了一个光锥微扰论中逐图进行的论证，其中不出现次序权重\cite[第~4~节]{Diehl:2015bca}（与\cref{app:appO}框架的比较见\cref{rem:appO-literature-comparison}）。余项的幂次：重整化子分析表明，重求和的 DY 截面中领头的红外重整化子相消，高扭度修正预期为 $\Lambda^2/[Q^2(1-z)^2]$ 量级（$Q(1-z)>\Lambda$ 时）\cite{BenekeBraun1995b}，这是证据而不是证明；$1/Q^2$ 处的广义因子化是\cref{thm:dy-higher-twist}，其外的幂次上因子化被破坏（\cref{prop:dy-twist6-failure}）。要把本定理提升到条件链完全成立的形式，需要证明 (3) 中的各项；要把它变成关于强子的陈述，还需要\cref{asm:intro-bridge}。



\cref{tab:dy-status}中，“在假设下完整”指该步在所列编号假设之下有完整证明；Glauber 特有的未证明输入是\cref{asm:appO-mixing}(c)、\cref{asm:appO-deformation}(a)--(d)、一般软子图的\cref{asm:appO-mixing}(a$'$)(b)、\cref{asm:appO-envelope}与\cref{asm:appO-rewritten-cs}，而附录~N 的输入与 DIS、$e^+e^-$ 共享。

\begin{remark}[固定阶部分子截面不检验旁观者 Glauber 相消]\label{rem:dy-nnlo-not-glauber}
DY 的部分子截面已经算到 NNLO\cite{HambergVanNeervenMatsuura1991}与 N$^3$LO（中性流\cite{Duhr:2020seh}与带电流\cite{Duhr:2020sdp}），质量因子化之后都是有限的，并与数据相符。但部分子截面里没有强子的旁观者：其有限性检验的是活动部分子的 Glauber（Coulomb）贡献在该阶的相消——例如单圈形状因子的虚部，即活动 $q\bar q$ 对之间红外发散的 Coulomb 相位，在与 Born 振幅的干涉 $2\Real F$ 中消去（\cref{sec:dy-nlo}）——而不检验旁观者之间的交换；后者是\cref{prop:dy-abelian-model}与\cref{app:appO}的内容。所以高阶计算与数据的符合是对整个框架的证据，而不是对旁观者 Glauber 相消的证明；“NNLO/N$^3$LO 的 DY 与数据一致，说明 Glauber 相消已被检验”是一个误解。
\end{remark}

\begin{remark}[Doria--Frenkel--Taylor 非相消的前提]\label{rem:dy-dft-premises}
Doria、Frenkel 与 Taylor 给出了一个反应，其中两个有质量、颜色平均的初态夸克的两软胶子非领头红外发散在实、虚图之间不相消\cite{DoriaFrenkelTaylor1980}；这一结论也被\cite{DiLietoEtAl1981}（他们在若干中间结果上与前者不一致）与\cite{Andrasi:1980qw}得到。这些非相消项在高能下被幂次压低：Frenkel、Gatheral 与 Taylor 表明到 $\order{g^4}$ 为止除了幂次压低的项之外是有限的\cite{Frenkel:1983di}，Catani 等人称之为“higher twist type”\cite{Catani:1985xt}（以上均据各文摘要）；对 DY，CSS 指出两圈的非相消出现在相对 $m^4/s^2$ 的量级\cite[第~10.2~节]{CollinsSoperSterman1989}。这一情形在\cref{thm:appO-main}的前提之外：后者处理由色单态顶点产生的强子（其旁观者携带补偿的色荷，并可与软胶子耦合）中无质量的活动部分子，并且只断言领头幂次；DFT 的外态则是有质量、有色的入射夸克，对其颜色求平均并不等价于色单态强子。在有质量的入射线上，$p_A=(p_A^+,\frac{m^2}{2p_A^+},\mathbf 0_T)$，
\begin{equation}\label{eq:dy-dft}
(p_A-k)^2-m^2=-2p_A^+\Big(k^-+\frac{m^2}{2(p_A^+)^2}\,k^+\Big)+k^2 ,
\end{equation}
% VF-REF-8
其中 $\frac{m^2}{2(p_A^+)^2}k^+$ 一项把 $k^+$ 与 $k^-$ 绑在一起，破坏了“$J_A$ 只通过 $k^-$ 感知软动量”这一使两个喷注可以分别按各自的光锥时间排序的性质（\cref{lem:dy-poles-orientation}所依据的领头幂次近似）；它是幂次压低的。对强子的 DY，因子化在 $m^4/s^2$ 处失效：在标量夸克模型中，因子化在 $m^2/s$ 成立而在 $m^4/s^2$ 失效\cite{Basu:1984ba}（\cref{prop:dy-twist6-failure}）。比较\cref{prop:ir-bn-failure}（\tagV）。
\end{remark}
